\section{The low sensitivity case}
\label{sur:section}

The goal in this section is to turn a small change in optimum cost
under budget reduction into a strict improvement below $2$.  We start
from the surplus solution, improve it by local search, and use its
same replay to bound the opening price.  If the bad clusters carry
little cost, the weighted swap bound suffices.  Otherwise enough of
their cost has bounded price for strictness to improve the total while
preserving every regular-facility payment.

We now fix the parameters used for this part of the algorithm.  Let
$C_0=C_{1/100}$ be the constant in Theorem~\ref{loc:anchoring}; explicitly,
the construction there permits
\[
 C_0=10+60000\,2^{3999}.
\]
Set $T=200C_0$ and $c=c(T)$ as in Lemma~\ref{str:strictness}.
Choose a fixed positive rational $v$ in the range of
Theorem~\ref{imp:cglss} satisfying
\begin{equation}\label{sur:parameter-choices}
 v\le\min\left\{\frac1{12},\frac{c}{1000}\right\},
 \qquad v^2\le\frac{c}{6(T+2)},
 \qquad \delta=3v.
\end{equation}
Fix $0<\xi\le1$; the final assembly will choose it smaller still.
Choose a fixed integer $N_*$ so large that for every $N\ge N_*$,
\begin{equation}\label{sur:size-threshold}
 N^2 2^{-N}\le v,\qquad 2/N\le c/200.
\end{equation}
Such an $N_*$ exists because $N^2 2^{-N}\to0$ and $2/N\to0$.
Enlarge the common integer surplus bound to satisfy
$L\ge\max\{1,N_*\}$ as well as the bound of
Theorem~\ref{imp:cglss}.  It still has order
$O_v(\log(N+2))$, with all preceding constants fixed.  In particular,
an instance with $k>2L$ has $N\ge k>2N_*$ and obeys
Equation~\eqref{sur:size-threshold}.  Instances of the finitely many
smaller sizes fall into the small-cardinality branch of the final
assembly.

\begin{proposition}[Low sensitivity]\label{sur:low-sensitivity}
With these fixed parameters, suppose $k>2L$ and
\begin{equation}\label{sur:sensitivity}
 \operatorname{opt}_{k-2L}\le(1+\xi)\operatorname{opt}_k.
\end{equation}
Put $h=k-L$.  In polynomial time one can compute a real facility set
$S$ of size $k$ with
\begin{equation}\label{sur:conclusion}
 \operatorname{cost}(S)
 \le\left(2-\frac{c}{100}\right)\operatorname{opt}_h
 \le\left(2-\frac{c}{100}\right)(1+\xi)
       \operatorname{opt}_k.
\end{equation}
The algorithm does not need to test Equation~\eqref{sur:sensitivity}.
\end{proposition}

\begin{proof}
Since $L\ge1$ and $k>2L$, we have $h\ge2$ and $h-L=k-2L\ge1$.
Apply Theorem~\ref{imp:cglss} at target $h$, then replay its output
using Lemma~\ref{imp:surplus-history} to obtain the real set $U$ and
its compatible budgets. In either output branch, $|U|\le h+L=k$.
Pad $U$ to size $k$ using members of $F$, and run width-$5$ local
search starting there; call its output $S$.  Padding is possible because
$k\le|F|$, and
\begin{equation}\label{sur:local-start}
 |S|=k,\qquad \operatorname{cost}(S)\le\operatorname{cost}(U).
\end{equation}
This is polynomial time on the bounded integral metric, as in
Lemma~\ref{loc:descent}.  Fix an optimum $O$ of size $h$, padding if
necessary, and write $P=\operatorname{opt}_h>0$,
$C_i=\{p:p\text{ is assigned to }i\in O\}$,
$n_i=|C_i|$, and $P_i=\sum_{p\in C_i}d(i,p)$.

\textit{The unmerged small-price branch.}
Suppose the output is the unmerged solution at
$\lambda=2/N^2$, with at most $h$ facilities.  Its payment guarantee is
\[
 (1-\delta)\operatorname{cost}(U)\le\sum_{p\in D}\alpha_p.
\]
Every nonempty cluster has $P_i\ge1$ in the positive integral metric,
so $\lambda\le2\le TP_i$.  Apply Lemma~\ref{str:strictness} to all
nonempty clusters.  Empty clusters contribute no budgets.  There are
at most $h$ nonempty clusters, and therefore
\begin{align*}
 (1-\delta)\operatorname{cost}(U)
 &\le h\lambda+(2-c/2)P\\
 &\le\left(2-c/2+c/200\right)P.
\end{align*}
Here $h\lambda\le2/N\le c/200\le(c/200)P$, because $h\le N$ and
$P\ge1$.  Finally
\[
 \frac{c}{200}+(2-c/100)\delta
 \le\frac{c}{200}+6v
 \le\frac{11c}{1000}<\frac{49c}{100}
 =\frac c2-\frac c{100}.
\]
It follows that $\operatorname{cost}(U)\le(2-c/100)P$, and
Equation~\eqref{sur:local-start} proves the desired conclusion in this
branch.  This argument uses only the payment bound with its opening
term dropped and the strictness just proved.

\textit{Dual comparison with an arbitrary real solution.}
It remains to consider the normal branch, whose replay has exactly $h$
paid regular facilities.  Put $e=hN\eta$, where $\eta=2^{-N}$.
Equation~\eqref{sur:size-threshold} ensures $e\le v\le vP$.
The payment and dual inequalities of Lemma~\ref{imp:surplus-history}
are
\begin{equation}\label{sur:payment-dual}
 (1-\delta)\operatorname{cost}(U)+h\lambda-e
 \le\sum_p\alpha_p,
 \qquad \sum_p[\alpha_p-2d(i,p)]_+\le\lambda\quad(i\in F).
\end{equation}
For any nonempty real set $H\subseteq F$, assign clients to nearest
centers of $H$.  For its cluster at $i$, the inequality
$\alpha_p\le2d(i,p)+[\alpha_p-2d(i,p)]_+$ gives, upon summing,
\begin{equation}\label{sur:any-H}
 \sum_p\alpha_p\le |H|\lambda+2\operatorname{cost}(H).
\end{equation}
Combining Equations~\eqref{sur:payment-dual} and~\eqref{sur:any-H}
and dropping the nonnegative connection term gives
\begin{equation}\label{sur:price-comparison}
 (h-|H|)\lambda\le2\operatorname{cost}(H)+e.
\end{equation}
Use for $H$ an optimum with $h-L=k-2L$ centers.  Equation~\eqref{sur:sensitivity}
and monotonicity of the optimum imply
$\operatorname{cost}(H)\le(1+\xi)\operatorname{opt}_k\le(1+\xi)P$,
and hence
\begin{equation}\label{sur:surplus-price}
 L\lambda\le2(1+\xi)P+e.
\end{equation}

\textit{Short clusters and simultaneous removals.}
Set $b=\lambda/8$, and for $i\in O$ let
$\operatorname{sep}_i=\min_{j\in O\setminus\{i\}}d(i,j)$.
Call $i$ short if $n_i\operatorname{sep}_i<b$; let $N_{\rm short}$
be the number of short centers.  Fix a nearest other center $\nu(i)$
for every $i\in O$.  Color all centers independently with two colors,
and select the short centers of the first color whose designated
neighbor $\nu(i)$ has the second color.  Each short center is selected
with probability $1/4$, so some coloring gives a set $R$ with
\[
 N':=|R|\ge N_{\rm short}/4.
\]
For every $i\in R$, $\nu(i)\notin R$.  Thus $O\setminus R$ is
nonempty, and all clients of $C_i$ can be rerouted to $\nu(i)$ with
total increase at most $n_i\operatorname{sep}_i<b$.  Keeping all other
assignments proves
\[
 \operatorname{cost}(O\setminus R)\le P+bN'.
\]
Empty clusters cause no difficulty: their rerouting increase is zero.
Apply Equation~\eqref{sur:price-comparison} to $H=O\setminus R$:
\[
 \lambda N'\le2P+2bN'+e.
\]
As $2b=\lambda/4$, we obtain
\begin{equation}\label{sur:short-price}
 \lambda N_{\rm short}\le4\lambda N'
 \le\frac{32}{3}P+\frac{16}{3}e.
\end{equation}
Equations~\eqref{sur:surplus-price} and~\eqref{sur:short-price} now
show explicitly that
\begin{align}
 b(L+N_{\rm short})
 &\le\left(\frac{19}{12}+\frac\xi4\right)P+\frac{19}{24}e
 \le10P.
 \label{sur:anchoring-scale}
\end{align}
For the last inequality use $\xi\le1$, $e\le vP$, and $v\le1/12$.
Since $|S|-h=L$, Equation~\eqref{sur:anchoring-scale} is precisely the
scale hypothesis of Theorem~\ref{loc:anchoring} with tolerance $1/100$.

\textit{The bad-cluster count and cost dichotomy.}
Apply Theorem~\ref{loc:anchoring} to $S,O,b$.  We obtain distinct proxies in $S$ for
the good centers, with total proxy assignment cost at most $P_g+P/100$,
and $m$ bad centers satisfying
\begin{equation}\label{sur:bad-price}
 m\lambda=8mb\le8C_0P.
\end{equation}
Write $P_b=\sum_{i\text{ bad}}P_i$ and $P_g=P-P_b$.
If $P_b\le P/10$, Lemma~\ref{loc:swap}, with $u=6/5$, gives
\[
 \operatorname{cost}(S)
 \le\frac65(P_g+P/100)+\frac{17}{5}P_b
 \le1.432P<\left(2-\frac c{100}\right)P,
\]
where $c<1$ follows from Equation~\eqref{str:constants}.

Otherwise $P_b>P/10$.  Let
\[
 \mathcal B_T=\{i\in O:i\text{ bad and }\lambda\le TP_i\},
 \qquad Q=\sum_{i\in\mathcal B_T}P_i.
\]
For every bad center excluded from $\mathcal B_T$, $P_i<\lambda/T$.
Equation~\eqref{sur:bad-price} and $T=200C_0$ therefore imply
\[
 \sum_{\substack{i\text{ bad}\\i\notin\mathcal B_T}}P_i
 \le m\lambda/T\le P/25,
 \qquad Q>P/10-P/25=3P/50>P/20.
\]
Every member of $\mathcal B_T$ has a nonempty positive-cost cluster,
so Lemma~\ref{str:strictness} applies there.  For each of the other
centers, the cluster version of Equation~\eqref{sur:any-H} gives
$\sum_{p\in C_i}\alpha_p\le\lambda+2P_i$, including empty clusters.
Summing over all $h$ centers of $O$ proves
\begin{equation}\label{sur:summed-gain}
 \sum_p\alpha_p
 \le h\lambda+2P-\frac c2Q
 \le h\lambda+\left(2-\frac c{40}\right)P.
\end{equation}
The full $h\lambda$ in Equation~\eqref{sur:summed-gain} cancels the
full $h\lambda$ in Equation~\eqref{sur:payment-dual}, giving
\[
 (1-\delta)\operatorname{cost}(U)
 \le\left(2-\frac c{40}\right)P+e.
\]
To divide by $1-\delta$, note that
\[
 \frac eP+(2-c/100)\delta
 \le v+6v\le\frac{7c}{1000}<\frac{3c}{200}
 =\frac c{40}-\frac c{100}.
\]
Thus $\operatorname{cost}(U)\le(2-c/100)P$, and
Equation~\eqref{sur:local-start} proves the first inequality in
Equation~\eqref{sur:conclusion}.  Finally
$P=\operatorname{opt}_h\le\operatorname{opt}_{h-L}
\le(1+\xi)\operatorname{opt}_k$ proves its second inequality.
All choices made by the algorithm were the imported construction,
padding, and local search; $O$, the coloring, and the good/bad partition
are used only in the proof.
\end{proof}
